\section{Asynchronous Inferential Cascade: A $\to$ B $\to$ C}
\label{sec:cascade-abc}

This section specifies a three-stage asynchronous inferential cascade layered
on top of the existing \textsc{Probingnoise} core. The cascade is a proposed
architectural extension and is not empirically validated within the present
work. The \textsc{Probingnoise} core remains the implemented synchronous,
Colab-trainable base; its blocks, residual pathways, and training regime are
preserved.

\subsection{Positioning}

The three stages address three distinct regimes of asynchrony identified as
architecturally independent.

\paragraph{Stage A --- multi-rate ingestion.} Heterogeneous input streams
(e.g., visual frames, auxiliary sensors, discrete events) arrive at different
timescales. Stage A normalizes them into a common representational substrate
without imposing a global synchronization barrier. Its architectural role is
to decouple upstream rate from downstream computation. Prior work on
event-driven and asynchronous graph processing~\citep{faber2022async,faber2024gwac}
establishes that useful representation can be maintained without strict
message-level synchronization; we adopt this only as lineage.

\paragraph{Stage B --- heterogeneous compute.} The \textsc{Probingnoise} core
operates within Stage B as the spatiotemporal backbone. Sub-processes within
this stage have unequal computational cost, and a synchronous schedule would
block fast sub-processes behind slow ones. Stage B is specified to permit
non-blocking progress across sub-processes, drawing architectural precedent
from lock-free stochastic optimization~\citep{niu2011hogwild} and foundational
work on asynchronous iterative methods~\citep{bertsekas1989parallel}. No
property of those works is inherited as a property of the current
implementation.

\paragraph{Stage C --- anytime refinement.} The topmost stage is specified to
emit partial predictions and continue refining in the background, rather than
blocking output on convergence. This design follows the general
anytime-algorithm regime; the contribution here is integration with Stages A
and B within a consistent cascade, not the anytime principle itself.

\subsection{Contribution boundary}

Novelty is claimed at the level of the integrated cascade --- specifically,
the assignment of one asynchronous regime to each stage and the specification
of their composition on top of the \textsc{Probingnoise} residual structure
--- rather than at the level of any individual stage component. Empirical
validation of the cascade is deferred to future work.

\subsection{Invariants}

\begin{itemize}
  \item Downstream stages never block on strict ordering guarantees from upstream stages.
  \item Each stage has an explicit definition of what ``partial output'' means for it.
  \item The \textsc{Probingnoise} core is preserved as-is within Stage B.
\end{itemize}

\subsection{Composition rule}

$A \to B \to C$ is a directional cascade. Each boundary is an async buffer,
not a synchronous hand-off. Backpressure is handled per-stage (drop-oldest by
default).
